% =============================================================================
% Project : THE_PHILOSOPHY_OF_VIABILITY
% Book    : Book I — Ontology of Viability
% Book ID : VS.B01
%
% File    :
%   2_MAINMATTER/PART_I_IDENTITY_AND_MANIFESTATION/
%   CHAPTER_02_INVARIANCE/sections/
%   section_06_invariance_and_structure.tex
%
% Section : 06 — Invariance and Structure
% Chapter : 02 — Invariance
%
% Purpose :
%   To examine structural invariance as a specified form of
%   preservation, distinguishing the preservation of components,
%   relations, and organizational arrangements from the numerical
%   identity of the existent in which they obtain.
%
% Conceptual Scope :
%   Structure, structural determination, component variation,
%   relational organization, structural transformation,
%   structural equivalence, and the limits of structural explanation.
%
% =============================================================================

\section{Invariance and Structure}
\label{sec:invariance-and-structure}

\subsection{The Ontological Relation Between Structure and Invariance}

Structure and invariance are related because an existent
may retain a specified structural determination across
a transformation. Changes in components, relations,
configurations, or material constitution do not necessarily
alter every aspect of the organization under consideration.

The preservation of a specified structural determination
under a defined transformation constitutes a particular
form of invariance, which may be described as structural
invariance.

Structural invariance concerns the preservation of an
identified organizational feature under specified
conditions. It does not refer simply to the continued
existence of the subject, nor does it establish that
the existent remains numerically identical through
every transformation under consideration.

Structure and identity address different ontological
questions. Structure concerns an organization or
relational determination of an existent. Identity
concerns the existent as such and cannot be reduced
to the preservation of any particular organizational
arrangement.

A structural determination may remain preserved
across a transformation without establishing the
numerical identity of the subjects before and after
that transformation. Conversely, an existent may
undergo structural alteration without that alteration
alone establishing the loss of its identity.

The distinction is therefore not between structure
and an existent considered as though they were
independent objects. It is between a particular
determination of an existent and the ontological
question of the identity of that existent.

Structural invariance must consequently be examined
as a determinate form of invariance, not as a general
explanation or criterion of identity.

\subsection{Structure as an Organized Determination}

A structure is not adequately characterized as
a mere aggregate of components.

An aggregate may be considered in terms of the
elements it contains, whereas structural determination
concerns the organization and relations through
which components or other determinations constitute
an arrangement.

The structural character of an existent may involve:

\begin{itemize}
    \item the organization of its components;
    \item the relations connecting those components;
    \item the arrangement of relations within a
          specified configuration;
    \item the constraints governing admissible
          configurations;
    \item the dependencies through which particular
          organizational determinations are maintained.
\end{itemize}

These aspects are analytically distinguishable.
They need not remain equally preserved under
a transformation.

A transformation may alter particular components
while preserving some relations among them. It
may preserve a relational arrangement while
altering other organizational constraints.
It may also modify the configuration as a whole
while retaining a more restricted structural feature.

For this reason, the assertion that a structure
remains invariant requires specification.

It is insufficient to state that ``the structure
remains the same'' without identifying the
organizational determination being considered,
the transformation under examination, and the
criterion according to which preservation is
established.

Structural invariance is therefore relative
to a specified structural respect. It cannot
be attributed to an entire structure without
clarifying which of its aspects are included
in the claim.

\subsection{Structural Preservation and Component Variation}

One possible form of structural invariance occurs
when components change while a specified
organizational relation remains preserved.

The elements occupying particular structural
positions may be modified or replaced without
necessarily altering the organization assessed
under a defined criterion.

For example, two relational configurations may
contain different elements while preserving
a corresponding arrangement of relations.
In such a case, the preservation concerns
the specified relational organization, not
the numerical identity of the components.

This possibility demonstrates that structural
invariance cannot be reduced to the persistence
of individual components.

However, the converse must also be recognized.

The persistence of components does not guarantee
the preservation of their structural organization.
The same components may enter into different
relations, occupy different positions, or become
organized according to a different arrangement.

Component preservation and structural preservation
are therefore analytically distinct.

Neither can be inferred automatically from the
other.

Furthermore, the replacement of a component
does not necessarily preserve the relevant
structure. Whether the organization remains
preserved depends on the structural respect
and criterion under consideration.

The fact that different components occupy
corresponding positions is insufficient unless
the relevant relations and organizational
conditions are also specified.

Structural invariance must consequently be
established through the assessment of the
organization itself, rather than inferred
solely from the persistence or replacement
of its components.

\subsection{Structural Relations and Organizational Preservation}

The preservation of a structure may concern
more than the persistence of individual
relations considered separately.

An organization may depend on how several
relations are connected, ordered, or mutually
constrained.

Consequently, the preservation of individual
relations does not necessarily establish
the preservation of their overall organization.

A configuration may retain several relations
while altering the arrangement through which
those relations constitute a structure.

Conversely, a specified organizational
arrangement may remain preserved even when
some individual relations are realized through
different components or determinations.

The distinction requires separating at least
three possible subjects of structural assessment:

\begin{enumerate}
    \item \textbf{Individual relations:}
    Whether particular relations remain preserved
    under the transformation.

    \item \textbf{Relational organization:}
    Whether the arrangement and interdependence
    of relations remain preserved according
    to a specified criterion.

    \item \textbf{Structural constraints:}
    Whether the conditions restricting or
    governing admissible configurations remain
    preserved.
\end{enumerate}

These subjects may be related, but they are
not interchangeable.

The preservation of individual relations
does not automatically establish the
preservation of their organization. Likewise,
the preservation of a structural constraint
does not necessarily establish that every
relation governed by that constraint remains
unchanged.

A structural invariance claim must therefore
identify whether it concerns individual
relations, their organization, particular
constraints, or a specified combination
of these features.

Without this specification, the assertion
of structural preservation remains
insufficiently determined.

\subsection{Structural Transformation and the Limits of Preservation}

Structural transformations may occur at
different analytical levels and with
different consequences for the organization
under consideration.

A transformation may:

\begin{enumerate}
    \item alter individual components while
          preserving a specified organizational
          relation;

    \item alter some relations while preserving
          a broader structural arrangement;

    \item reorganize relations while retaining
          particular structural constraints;

    \item alter the organizational arrangement
          to such an extent that the previously
          specified structural feature is no
          longer preserved.
\end{enumerate}

These possibilities indicate that structural
change is not a single, uniform phenomenon.

A structure may remain invariant with respect
to one transformation and fail to remain
invariant with respect to another. It may
also preserve one structural respect while
undergoing alteration in another.

Thus, structural invariance cannot be asserted
without specifying the transformation or
class of transformations, the relevant domain,
the structural respect, and the criterion
of preservation.

The recognition of structural change also
does not establish, by itself, that the
identity of the existent has been lost.

Such a conclusion would require an independent
ontological argument concerning identity.
The alteration of an organizational arrangement
may be relevant to that investigation, but
it cannot independently settle it.

Likewise, the preservation of a structural
arrangement cannot establish that the same
existent has persisted. A corresponding
or equivalent organization may obtain in
distinct existents.

Structural transformation and structural
preservation must therefore be examined
without presuming that either determines
the numerical identity of the existent.

\subsection{Structural Equivalence and Numerical Identity}

A particularly important distinction concerns
structural equivalence and numerical identity.

Two existents may exhibit corresponding
structural organizations without being
numerically identical.

Their structures may be equivalent under
a specified criterion, while the existents
in which those structures obtain remain
distinct.

Structural equivalence expresses a relation
between structural determinations. Numerical
identity concerns whether an existent is
one and the same existent.

These are different ontological questions.

The distinction prevents the reduction
of identity to structure.

If structural equivalence were treated
as sufficient for numerical identity,
distinct existents possessing corresponding
organizations would incorrectly be treated
as one existent.

The same problem arises when the preservation
of a structure through transformation is
taken as proof of the persistence of the
existent associated with it.

Such an inference exceeds what structural
invariance alone can establish.

Structural equivalence may support a
comparison between configurations, and
structural invariance may establish the
preservation of a specified organization.
Neither conclusion independently settles
the numerical identity of the existents
being compared.

The ontological status of the subjects
must therefore remain distinct from the
equivalence or invariance of their
structural determinations.

\subsection{Structural Sameness and the Limits of Structural Identity}

Expressions such as ``the same structure''
or ``structural identity'' require
qualification.

They may refer to the preservation of
an organizational arrangement across
a transformation, or to the correspondence
of two structures under a specified
criterion.

In either case, the expression concerns
the sameness or equivalence of structural
determinations. It must not be confused
with the ontological identity of the
existent whose structure is being examined.

A structure may be treated as the same
under a defined criterion even when
the components occupying its positions
have changed.

Likewise, two structures may be described
as equivalent under a formal criterion
while belonging to distinct existents.

Such descriptions establish a particular
analytical relation. They do not establish
that the associated existents are
numerically one and the same.

Moreover, an existent must not be assumed
to be exhausted by a single structural
determination or by one structural
description.

Different structural respects may be
relevant to the same existent. A
transformation may preserve one of
these respects while altering another.

Consequently, any assertion of structural
sameness must specify its subject,
criterion, and scope.

The expression cannot serve as an
unqualified substitute for ontological
identity, nor can a particular structural
description be presumed to exhaust
the determinations of the existent.

\subsection{The Limits of Structural Explanation of Identity}

The preceding analysis establishes
a necessary limitation: structural
invariance is a possible form of
preservation, but the identity of
an existent cannot be reduced to
the preservation of its structure.

This limitation operates in both
directions.

First, structural preservation does
not prove numerical identity.

Equivalent structures may belong
to distinct existents, and the
preservation of an organizational
arrangement does not, by itself,
establish the persistence of the
existent associated with it.

Second, structural transformation
does not prove the loss of identity.

An existent may undergo changes
in its organization without the
mere occurrence of those changes
determining that it has become
another existent.

Third, the absence of an established
structural invariant does not establish
the absence of identity.

The failure to identify a preserved
structural feature may indicate
that the relevant structure has
changed, that the criterion is
inadequate, or that the available
analysis does not determine the
matter. It cannot independently
settle the ontological question
of identity.

Structure may therefore contribute
to the analysis of an existent
without becoming an exhaustive
ground of its identity.

Identity and structure must remain
analytically distinct. The fact
that structure is a determination
of an existent does not make it
identical with the existent's
identity, nor does it establish
that structural preservation is
a necessary or sufficient condition
of numerical identity.

The proper conclusion is not that
structure is ontologically irrelevant
to identity, but that its relevance
must not be converted into reduction.

\subsection{Synthesis}

Structural invariance concerns the
preservation of a specified organizational
determination across a defined domain
of transformation.

Its analysis requires distinguishing:

\begin{itemize}
    \item the preservation of components
          from the preservation of their
          organization;

    \item the preservation of individual
          relations from the preservation
          of their relational arrangement;

    \item the preservation of structural
          constraints from the preservation
          of every configuration governed
          by those constraints;

    \item structural equivalence from
          numerical identity;

    \item structural transformation from
          the loss of ontological identity;

    \item structural sameness under a
          specified criterion from the
          identity of the existent itself.
\end{itemize}

The preservation of structure is neither
a sufficient condition for numerical
identity nor a universally necessary
condition that can be imposed upon
every existent.

Structural invariance remains an
important analytical category within
the ontology of invariance. Its validity,
however, is restricted to the structural
determination, transformation, domain,
and criterion of preservation under
consideration.

The identity of an existent must
consequently remain an independent
ontological question, irreducible
to the invariance, equivalence, or
transformation of its structure.

The next section examines the relation
between invariance and composition,
addressing the distinction between
the preservation of an organized
arrangement and the preservation
or alteration of the constituents
through which that arrangement
is constituted.

% =============================================================================
% END OF SECTION 06 OF CHAPTER 02
% =============================================================================